\newpage
\onecolumn
\appendix

\subsection{Contributions}
The project was started by HW and DD. DD, HW, and JTS designed the video encoder and short-horizon memory system. KP designed the long-horizon memory system. MT, KP, HW, KV, SN, and ME designed the evaluation suite and performed experimental evaluations, including model post-training. DD, KP, AR, and QV developed the pre-training recipe and data mixture, and ran large-scale model training. HW, KD, KS, and MT developed the infrastructure for long-context VLA training. JT optimized inference of the \Acronym{} VLA. SL and CF provided advice throughout the project. MT, KP, DD, BI, KV, SN, SL, and CF worked on writing, illustrations, and the video. 

\subsection{Task Details}
\label{sec:task_details}

\subsubsection{Long-horizon Tasks (\cref{sec:challenge_tasks})}

\paragraph{Preparing the items for a recipe} The goal of the task is for the robot to take out all of the items for a recipe. The items are placed in various, randomized locations throughout the kitchen (fridge, cabinet, drawers, stove ...). The policy's prompt specifies which items should be assembled, where they should be placed (countertop, stove, sink), and their respective initial location in the kitchen. We list a few prompt examples below. The robot will receive 1 point per item retrieved and placed in the requested position. Most recipes consist of 6 to 7 points. We collected data on over 40 recipes with a large variety of objects, scenes, appliances, prompts. The evaluation is done on a variety of unseen scenes, with seen recipes. 
\begin{PromptBox}{Mashed potatoes}
Take out supplies for the recipe. Remove the lid from the pot on the stove and place the lid on the countertop to the left of the sink, and pick up the pot and place it in the sink. Get the potatoes, butter, milk from the fridge and place them on the countertop to the left of the sink. Get the masher from the top drawer to the left of the sink and place it on the countertop to the left of the stove.
\end{PromptBox}

\begin{PromptBox}{Fried rice}
Take out the supplies for the recipe. Get the bag of rice and the spam from the cabinet on the bottom right of the dishwasher and place them on the countertop to the left of the sink. Get the soy sauce from the fridge and place it on the same section of the countertop to the left of the sink. Get the pan from the bottom cabinet to the right of the dishwasher and put it on the stovetop. Get the spatula from the top drawer to the right of the dishwasher and place it on the pan on the stovetop.
\end{PromptBox}

\begin{PromptBox}{Pizza}
Take out the supplies for the recipe. Get the baking tray from the bottom cabinet on the left of the dishwasher and place it on the countertop to the left of the sink. Get the pizza dough, package of pepperoni, and bag of cheese from the fridge and place them on the countertop to the left of the sink. Get the dough roller from the top drawer to the right of the stove and place it on the countertop to the left of the sink.
\end{PromptBox}
 

\paragraph{Cleaning the kitchen} This task requires cleaning a kitchen by wiping the countertop with the sponge, drying the countertop with a paper towel, storing any food items left in the countertop inside the fridge, putting all the dishes from the dishrack into the cabinets, and washing all of the dishes from the sink. The scoring consists of +1 point per subtask done and on average episodes consist of 8 subtasks. Below we provide an example prompt.
\begin{PromptBox}{Clean kitchen}
Clean up the kitchen. Wipe the countertop with the sponge and dry it with the paper towel and throw the towels in the trash can under the sink. Put the mustard in the fridge. Put the dishes in the bottom cabinet on the left of the fridge. Wash the blue plate and black plate and place them in the dishrack. 
\end{PromptBox}


\subsubsection{In-context adaptation (\cref{sec:in_context_exp})}

\paragraph{Chopstick Pick Up} The policy is tasked with picking up a chopstick on a variable height table. During data collection, the chopstick is placed at random locations on the table with random table heights in the upper half of the table height range. We collect human interventions whenever the policy mis-grasps, and then train both memory and non-memory policies on this data and evaluate them with the chopstick randomly placed on the table at the lowest table height setting. The final policies are scored +1 for picking up the chopstick and +1 for placing it in the bin, and success is a score of 2.

\paragraph{Open Refrigerator} The goal of the task is to open a fridge in front of the robot. The fridge does not have obvious visual cues indicating which side the door hinge is on, leading policies to often attempt to open it in the wrong direction. We collect human interventions whenever the robot opens the fridge in the wrong direction. For evaluation, an episode is defined as successful if it takes $\leq 4$ grasps to open the door, to test for intentional strategy switching rather than repeated random sampling.

\vspace{0.3cm}
\subsubsection{Analysis Experiments (\cref{sec:analysis})}
\paragraph{Three-way swap mugs} The goal is to place three coffee mugs under a coffee machine sequentially. Two mugs start at random positions on the table and a third mug starts under the coffee machine. The scoring consists of obtaining +1 point for each mug that goes on the coffee machine without tipping and without repeating mugs. Only one mug can go on the coffee machine at a time, so the robot needs to place it back on a random location on the table after it was placed under the coffee machine to make space for the next mug.

\paragraph{Find object} The goal is to retrieve the hidden object in a cabinet with four drawers. In the beginning of the episode, a person places the object in a random drawer. The robot needs to remember which drawer the person placed the object into, then open the correct drawer and retrieve the item and place it on the table. One point will be received for successfully accomplishing the task without opening any incorrect drawer.

\paragraph{Unpack groceries} The goal is to retrieve all items from a grocery bag without missing any items inside. The inside of the bag is not observable from the third person camera of the robot, and the number of items in the bag is randomized across episodes. Thus, to complete the task, the robot will need to remember how many objects are left in the bag from its past wrist camera observations. Episodes are marked as a success if the policy retrieved all items from the bag and indicated completion by moving home once all items are retrieved.

\paragraph{Scoop coffee} The goal is to put exactly two scoops of coffee beans in a grinder. The robot needs to remove the lid from the hopper of the grinder, pick up a coffee scoop, put two scoops of coffee from an opened coffee bean package into the grinder, then put the lid back on. Episodes are marked as a success if exactly two scoops are added to the grinder and the lid is put back on.

\paragraph{Grilled cheese} The goal of the task is to prepare a grilled cheese sandwich. The robot has to assemble the grilled cheese sandwich in the pan, then wait for the grilled cheese to cook, flip it, wait for the other side to cook, and finally plate the grilled cheese. The policy receives +1 point for assembling the sandwich, cooking for the correct amount of time (between 30 seconds and 3 minutes per side), flipping, and plating.

\paragraph{Window cleaning} The goal is to wipe the window door of a phone booth. The task consists of spraying windex on the window, ripping off a paper towel from a roll, then wiping the whole window, and finally throwing the paper towel in the trash. The robot needs to remember which steps have been completed and which part of the window have already been wiped during the cleaning phase. Episodes are marked as a success if all steps of the task are completed. It is considered a failure if the policy fails to wipe parts of the window.

\paragraph{Table bussing} The goal is to sort 12~objects from a tabletop into (A)~a bussing bin (for utensils and tableware), and (B)~a trash can (for plastic containers, bottles, and wrapping paper). The policy receives a score of +1 for every correctly placed item.

\paragraph{Shirt folding} The robot has to fold a shirt on a tabletop. At the beginning of the episode the shirt is flattened on the table. The episode is scored as a success if the shirt is successfully folded into a rectangle without seams sticking out or excessive wrinkles.

\paragraph{Clean up counter} The goal is to clean items from a counter into a drawer. The robot has to open the drawer, place all items on the counter into the drawer, and close it. An episode is counted as a success if all subtasks are completed successfully.

\paragraph{Make bed} The goal is to tidy a bed, by straightening the blanket, and placing any pillows at the head of the bed. The episode is scored as a success if all pillows are placed at the head of the bed, and the blanket is tidied.

\paragraph{Kitchen cleanup} The goal is to clean the counter of a kitchen, by placing multiple plates, bowls, and a cutting board from the counter into a nearby sink. The episode is scored as a success if all plates are placed into the sink.

\paragraph{Batch folding} In contrast to the shirt folding task, here the clothes start in a hamper on one side of the table. The robot needs to take the clothing item out of the hamper, flatten it, fold it, and then stack it onto a stack of existing clothes on the table. The hamper may contain various t-shirts and shorts. Episodes are scored as successes if the clothing items are folded into rectangles without seams sticking out or excessive wrinkling.

\paragraph{Box building} Given a flattened cardboard cutout, the robot is tasked to fold the cardboard into a box. This task requires multiple precise, well-coordinated bimanual manipulations to assemble the box. Episodes are scored as successful if the box is fully assembled without damaging the box.







\subsection{Video encoder with Space-Time separable attention}
\label{appendix:st_attention}
We describe how we can adjust a layer in a given ViT image encoder to instantiate our video-encoding scheme.

Let $\mathbf{z}^{l-1}_{p, t}$ denote the input embeddings to layer $l$ for spatial patch $p$ and timestep t (where $t \in [-K, 0]$). We first modify the embedding by adding a sinusoidal position embedding based on $t$, denote the output of this step with 
$$
\hat{\mathbf{z}}^{l-1}_{p, t} = \mathbf{z}^{l-1}_{p, t} + e(t),
$$
where $e(t)$ denotes the sinusoidal position embedding and we set the boundary condition $e(0) = 0$.

We then re-use the ViT's standard query key and value projections which in our case are given as
\begin{equation}
    \begin{aligned}
        \mathbf{q}^{l, a}_{p, t} &= W^{l, a}_Q \text{LN}\big( \hat{\mathbf{z}}^{l-1}_{p, t} \big), \\
        \mathbf{k}^{l, a}_{p, t} &= W^{l, a}_K \text{LN}\big( \hat{\mathbf{z}}^{l-1}_{p, t} \big), \\
        \mathbf{v}^{l, a}_{p, t} &= W^{l, a}_V \text{LN}\big( \hat{\mathbf{z}}^{l-1}_{p, t} \big),
    \end{aligned}
\end{equation}
where $a$ indexes the attention head and LN denotes a layer-norm implementation (we use RMSNorm but this choice is dependent on the ViT we start from).

Next, we can define the general attention mechanism (ignoring normalization for brevity of notation) as the softmax over the query and key outer product:
\begin{equation}
    \mathbf{\alpha}^{l,a}_{p,t} (\mathcal{S} = \lbrace1, \dots N \rbrace, \mathcal{T} = \lbrace 1, \dots, T \rbrace) [\hat{\mathbf{z}}] = \text{SM}\Big( (\mathbf{q}^{l, a}_{p, t})^T \cdot \big[ \mathbf{k}^{l, a}_{0, 0} \lbrace \mathbf{k}^{l, a}_{p', t'} \rbrace_{p' \in \mathcal{S}, t' \in \mathcal{T}} \big]\Big),
\end{equation}
where $\text{SM}$ denotes the softmax operation and $\mathcal{S} = \lbrace1, \dots N \rbrace$ and $\mathcal{T} = \lbrace 1, \dots, T \rbrace$ denote the space and time indices we want to perform attention over respectively.
With this general definition we can define the space only attention mechanism as instantiating $\mathbf{\alpha}^{l,a}_{p,t} (\mathcal{S} = \lbrace 1, \dots N \rbrace, \mathcal{T} = \lbrace \rbrace)$ and the time only attention mechanism as instantiating $\mathbf{\alpha}^{l,a}_{p,t} (\mathcal{S} = \lbrace \rbrace, \mathcal{T} = \lbrace 1, \dots, T \rbrace)$.
Thus the attention mechanism used in our video encoder is precisely given as
\begin{equation}
    \mathbf{\alpha}^{l,a}_{p,t} (\mathcal{S} = \lbrace 1, \dots N \rbrace, \mathcal{T} = \lbrace \rbrace) [ \mathbf{\alpha}^{l,a}_{p,t} (\mathcal{S} = \lbrace1, \dots N \rbrace, \mathcal{T} = \lbrace 1, \dots, T \rbrace)[\hat{\mathbf{z}}]].
\end{equation}
And thereafter we follow the standard computation of a transformer layer, i.e. we compute the outputs of the transformer by first combining attention weights $\alpha$ with values $\mathbf{v}$ and passing the corresponding outputs to a MLP.
